\textbf{Phần 1. Trắc nghiệm lựa chọn}
\setcounterpageref{ex}{0}

\begin{ex} %[4D1V04]
Cho biểu thức $A = 125 \times m + 875 \times m$. Với $m = 8$, giá trị của biểu thức $A$ là
\choice
{\True $8000$}
{$7500$}
{$8500$}
{$10\,000$}
\loigiai{
Áp dụng tính chất nhân một số với một tổng:
$$A = m \times (125 + 875) = m \times 1000$$
Với $m = 8$, ta có:
$$A = 8 \times 1000 = 8000$$
\textbf{Đáp án của câu hỏi là A}
}
\end{ex}

\begin{ex} %[4D1V04]
Một hình chữ nhật có chiều dài là $a$, chiều rộng là $b$ ($a, b$ cùng đơn vị đo xăng-ti-mét). Nếu tăng chiều dài thêm $5\text{ cm}$ và giữ nguyên chiều rộng thì chu vi của hình chữ nhật tăng thêm bao nhiêu xăng-ti-mét?
\choice
{$5\text{ cm}$}
{$15\text{ cm}$}
{$20\text{ cm}$}
{\True $10\text{ cm}$}
\loigiai{
Chu vi ban đầu của hình chữ nhật là: $P = (a + b) \times 2 = 2 \times a + 2 \times b$.\\
Khi tăng chiều dài thêm $5\text{ cm}$, chu vi mới là:
$$P_{\text{mới}} = (a + 5 + b) \times 2 = 2 \times a + 2 \times b + 10 = P + 10\text{ (cm)}$$
Vậy chu vi của hình chữ nhật tăng thêm $10\text{ cm}$.\\
\textbf{Đáp án của câu hỏi là D}
}
\end{ex}

\begin{ex} %[4D1V04]
Tìm số tự nhiên $x$ để biểu thức $Q = 20\,000 : (x - 3)$ đạt giá trị lớn nhất.
\choice
{$x = 3$}
{\True $x = 4$}
{$x = 5$}
{$x = 1$}
\loigiai{
Để phân số hoặc phép chia có số bị chia cố định đạt giá trị lớn nhất thì số chia phải là số tự nhiên bé nhất khác 0.\\
Do đó, ta có:
\begin{align*}
x - 3 &= 1 \\
x &= 1 + 3 \\
x &= 4
\end{align*}
Khi đó $Q = 20\,000 : 1 = 20\,000$ (giá trị lớn nhất).\\
\textbf{Đáp án của câu hỏi là B}
}
\end{ex}

\begin{ex} %[4D1V04]
Cho biểu thức $B = a \times 4 + b \times 4 + c \times 4$. Biết trung bình cộng của ba số $a, b, c$ bằng $25$. Giá trị của biểu thức $B$ là
\choice
{$100$}
{$200$}
{\True $300$}
{$400$}
\loigiai{
Vì trung bình cộng của $a, b, c$ bằng $25$ nên tổng của ba số là:
$$a + b + c = 25 \times 3 = 75$$
Ta có:
$$B = 4 \times (a + b + c) = 4 \times 75 = 300$$
\textbf{Đáp án của câu hỏi là C}
}
\end{ex}

\begin{ex} %[4D1V04]
Một hình chữ nhật có chu vi $P = 96\text{ cm}$. Nếu biết chiều dài là $a = 32\text{ cm}$ thì chiều rộng $b$ của hình chữ nhật đó là
\choice
{\True $16\text{ cm}$}
{$18\text{ cm}$}
{$32\text{ cm}$}
{$64\text{ cm}$}
\loigiai{
Nửa chu vi của hình chữ nhật là:
$$a + b = 96 : 2 = 48\text{ (cm)}$$
Chiều rộng $b$ của hình chữ nhật là:
$$b = 48 - 32 = 16\text{ (cm)}$$
\textbf{Đáp án của câu hỏi là A}
}
\end{ex}

\begin{ex} %[4D1V04]
Rút gọn biểu thức $M = (a + b) \times c - (a - b) \times c$, ta được biểu thức nào dưới đây?
\choice
{$2 \times a \times c$}
{$a \times b \times c$}
{\True $2 \times b \times c$}
{$0$}
\loigiai{
Áp dụng tính chất phân phối của phép nhân đối với phép trừ:
\begin{align*}
M &= c \times [(a + b) - (a - b)] \\
&= c \times [a + b - a + b] \\
&= c \times (2 \times b) = 2 \times b \times c
\end{align*}
\textbf{Đáp án của câu hỏi là C}
}
\end{ex}

\begin{ex} %[4D1V04]
Một đội công nhân làm đường trong 3 ngày. Ngày thứ nhất làm được $m$ mét, ngày thứ hai làm được nhiều hơn ngày thứ nhất $20$ mét, ngày thứ ba làm được gấp đôi ngày thứ hai. Biểu thức biểu thị tổng số mét đường đội công nhân làm được trong cả ba ngày là
\choice
{$3 \times m + 40$}
{\True $4 \times m + 60$}
{$4 \times m + 40$}
{$3 \times m + 60$}
\loigiai{
- Ngày thứ nhất làm được: $m\text{ (m)}$.\\
- Ngày thứ hai làm được: $m + 20\text{ (m)}$.\\
- Ngày thứ ba làm được: $2 \times (m + 20) = 2 \times m + 40\text{ (m)}$.\\
Tổng số mét đường làm được trong cả ba ngày là:
$$m + (m + 20) + (2 \times m + 40) = 4 \times m + 60\text{ (m)}$$
\textbf{Đáp án của câu hỏi là B}
}
\end{ex}

\begin{ex} %[4D1V04]
Cho biểu thức $H = 2024 - (a + b + c)$. Biết $a, b, c$ là ba số chẵn liên tiếp có tổng bằng $72$. Giá trị của biểu thức $H$ là
\choice
{$1942$}
{$1962$}
{$1950$}
{\True $1952$}
\loigiai{
Theo đề bài, tổng của ba số chẵn liên tiếp là $a + b + c = 72$.\\
Thay vào biểu thức $H$, ta được:
$$H = 2024 - 72 = 1952$$
\textbf{Đáp án của câu hỏi là D}
}
\end{ex}

\begin{ex} %[4D1V04]
Giá trị của biểu thức $K = m \times 16 + m \times 84 - 15\,000$ với $m = 250$ là
\choice
{\True $10\,000$}
{$15\,000$}
{$20\,000$}
{$25\,000$}
\loigiai{
Ta có:
\begin{align*}
K &= m \times (16 + 84) - 15\,000 \\
&= m \times 100 - 15\,000
\end{align*}
Với $m = 250$:
$$K = 250 \times 100 - 15\,000 = 25\,000 - 15\,000 = 10\,000$$
\textbf{Đáp án của câu hỏi là A}
}
\end{ex}

\begin{ex} %[4D1V04]
Một mảnh đất hình vuông có cạnh là $a$ (mét). Người ta mở rộng mảnh đất về một phía thêm $4\text{ m}$ để được mảnh đất hình chữ nhật. Biểu thức tính diện tích phần đất tăng thêm là
\choice
{$a \times a + 4$}
{$(a + 4) \times 4$}
{\True $a \times 4$}
{$a \times 2$}
\loigiai{
Phần đất tăng thêm là một hình chữ nhật có một kích thước là chiều dài cạnh hình vuông ban đầu ($a\text{ m}$) và một kích thước là phần mở rộng thêm ($4\text{ m}$).\\
Do đó, diện tích phần đất tăng thêm là:
$$S_{\text{tăng}} = a \times 4\text{ (m}^2\text{)}$$
\textbf{Đáp án của câu hỏi là C}
}
\end{ex}

\begin{ex} %[4D1V04]
Với $x = 120, y = 40, z = 5$, giá trị của biểu thức $(x - y) : z + x : z$ là
\choice
{$32$}
{\True $40$}
{$48$}
{$50$}
\loigiai{
Thay các giá trị vào biểu thức:
\begin{align*}
(120 - 40) : 5 + 120 : 5 &= 80 : 5 + 24 \\
&= 16 + 24 = 40
\end{align*}
\textbf{Đáp án của câu hỏi là B}
}
\end{ex}

\begin{ex} %[4D1V04]
Cho hai biểu thức $P = a \times 6 + b \times 6$ và $Q = (a + b) \times 5$. Nếu $a + b = 50$ thì giá trị của $P$ lớn hơn giá trị của $Q$ số đơn vị là
\choice
{$10$}
{$25$}
{$100$}
{\True $50$}
\loigiai{
Ta có:
$$P = 6 \times (a + b)$$
$$Q = 5 \times (a + b)$$
Hiệu giữa $P$ và $Q$ là:
$$P - Q = 6 \times (a + b) - 5 \times (a + b) = a + b$$
Vì $a + b = 50$ nên $P$ lớn hơn $Q$ đúng $50$ đơn vị.\\
\textbf{Đáp án của câu hỏi là D}
}
\end{ex}

\textbf{Phần 2. Trắc nghiệm đúng sai}
\setcounterpageref{ex}{0}

\begin{ex} %[4D1V04]
Xét tính đúng sai của các khẳng định sau về biểu thức chứa chữ:
\choiceTF
{\True Biểu thức $a + b - 15$ là biểu thức chứa hai chữ $a$ và $b$.}
{Với $m = 8$, giá trị của biểu thức $120 - 20 \times m$ bằng $800$.}
{\True Chu vi của hình tam giác có độ dài ba cạnh là $a, b, c$ (cùng đơn vị đo) được tính bởi biểu thức $P = a + b + c$.}
{\True Với $x = 15, y = 5$, giá trị của biểu thức $x \times y + x : y$ bằng $78$.}
\loigiai{
\begin{itemize}
    \item Biểu thức $a + b - 15$ chứa hai chữ cái đại diện cho hai biến số là $a$ và $b$. Khẳng định này Đúng.
    \item Với $m = 8$, ta thực hiện phép nhân trước: $20 \times 8 = 160$. Khi đó $120 - 160$ không thể thực hiện trong tập số tự nhiên (không bằng 800). Khẳng định này Sai.
    \item Chu vi tam giác bằng tổng độ dài 3 cạnh: $P = a + b + c$. Khẳng định này Đúng.
    \item Với $x = 15, y = 5$, giá trị biểu thức là: $15 \times 5 + 15 : 5 = 75 + 3 = 78$. Khẳng định này Đúng.
\end{itemize}
\begin{itemize}
    \item a) \textbf{Đúng}.
    \item b) \textbf{Sai}.
    \item c) \textbf{Đúng}.
    \item d) \textbf{Đúng}.
\end{itemize}
}
\end{ex}

\begin{ex} %[4D1V04]
Cho hình chữ nhật có chiều dài $a = 25\text{ cm}$ và chiều rộng $b = 15\text{ cm}$. Xét tính đúng sai của các khẳng định sau:
\choiceTF
{\True Chu vi của hình chữ nhật đó là $80\text{ cm}$.}
{\True Diện tích của hình chữ nhật đó là $375\text{ cm}^2$.}
{Nếu cùng tăng cả chiều dài và chiều rộng thêm $5\text{ cm}$ thì chu vi tăng thêm $10\text{ cm}$.}
{\True Nếu tăng chiều dài thêm $5\text{ cm}$ và giảm chiều rộng đi $5\text{ cm}$ thì chu vi hình chữ nhật không đổi.}
\loigiai{
\begin{itemize}
    \item Chu vi ban đầu là: $P = (25 + 15) \times 2 = 80\text{ (cm)}$. Khẳng định a Đúng.
    \item Diện tích ban đầu là: $S = 25 \times 15 = 375\text{ (cm}^2\text{)}$. Khẳng định b Đúng.
    \item Khi cùng tăng cả chiều dài và chiều rộng thêm $5\text{ cm}$ thì chu vi mới là: $((25 + 5) + (15 + 5)) \times 2 = (40 + 10) \times 2 = 100\text{ (cm)}$, tăng thêm $100 - 80 = 20\text{ (cm)}$ (không phải $10\text{ cm}$). Khẳng định c Sai.
    \item Khi chiều dài tăng $5\text{ cm}$ và chiều rộng giảm $5\text{ cm}$ thì tổng chiều dài và chiều rộng là: $(a + 5) + (b - 5) = a + b$, do đó chu vi không thay đổi. Khẳng định d Đúng.
\end{itemize}
\begin{itemize}
    \item a) \textbf{Đúng}.
    \item b) \textbf{Đúng}.
    \item c) \textbf{Sai}.
    \item d) \textbf{Đúng}.
\end{itemize}
}
\end{ex}

\textbf{Phần 3. Trả lời ngắn}
\setcounterpageref{ex}{0}

\begin{ex} %[4D1V04]
Cho biểu thức $A = 35 \times m + 65 \times m$. Tính giá trị của biểu thức $A$ với $m = 48$.
\shortans{4800}
\loigiai{
\textbf{Đáp số:} $4800$\\
Áp dụng tính chất phân phối:
$$A = (35 + 65) \times m = 100 \times m$$
Với $m = 48$:
$$A = 100 \times 48 = 4800$$
}
\end{ex}

\begin{ex} %[4D1V04]
Cho biểu thức $P = a + b + c$. Biết $a$ là số chẵn lớn nhất có hai chữ số, $b$ là số lẻ bé nhất có hai chữ số khác nhau, và $c = 45$. Tính giá trị của biểu thức $P$.
\shortans{156}
\loigiai{
\textbf{Đáp số:} $156$\\
- Số chẵn lớn nhất có hai chữ số là $a = 98$.\\
- Số lẻ bé nhất có hai chữ số khác nhau là $b = 13$.\\
Giá trị của biểu thức $P$ là:
$$P = 98 + 13 + 45 = 156$$
}
\end{ex}

\begin{ex} %[4D1V04]
Một hình chữ nhật có chu vi $P = (a + b) \times 2$. Biết chu vi $P = 120\text{ cm}$ và chiều dài $a = 38\text{ cm}$. Tính diện tích của hình chữ nhật đó theo đơn vị xăng-ti-mét vuông.
\shortans{836}
\loigiai{
\textbf{Đáp số:} $836$\\
Nửa chu vi của hình chữ nhật là:
$$a + b = 120 : 2 = 60\text{ (cm)}$$
Chiều rộng $b$ của hình chữ nhật là:
$$b = 60 - 38 = 22\text{ (cm)}$$
Diện tích của hình chữ nhật đó là:
$$S = a \times b = 38 \times 22 = 836\text{ (cm}^2\text{)}$$
}
\end{ex}

\begin{ex} %[4D1V04]
Cho biểu thức $M = x \times y - z$. Biết $x = 125$, $y$ là số chẵn lớn nhất có một chữ số và $z = 350$. Tính giá trị của biểu thức $M$.
\shortans{650}
\loigiai{
\textbf{Đáp số:} $650$\\
Số chẵn lớn nhất có một chữ số là $y = 8$.\\
Thay vào biểu thức $M$, ta có:
$$M = 125 \times 8 - 350 = 1000 - 350 = 650$$
}
\end{ex}

\begin{ex} %[4D1V04]
Cho $a + b = 100$ và $c + d = 50$. Tính giá trị của biểu thức:
$$T = a \times 5 + c \times 5 + b \times 5 + d \times 5$$
\shortans{750}
\loigiai{
\textbf{Đáp số:} $750$\\
Ta nhóm các số hạng:
\begin{align*}
T &= 5 \times (a + b + c + d) \\
&= 5 \times [(a + b) + (c + d)] \\
&= 5 \times [100 + 50] \\
&= 5 \times 150 = 750
\end{align*}
}
\end{ex}

\begin{ex} %[4D1V04]
Tìm số tự nhiên $a$ để biểu thức $Q = 15\,000 : (20 - a)$ có giá trị lớn nhất.
\shortans{19}
\loigiai{
\textbf{Đáp số:} $19$\\
Để biểu thức $Q$ đạt giá trị lớn nhất thì số chia $20 - a$ phải là số tự nhiên bé nhất khác 0.\\
Do đó:
\begin{align*}
20 - a &= 1 \\
a &= 20 - 1 \\
a &= 19
\end{align*}
}
\end{ex}
